% !TEX root = userguide.tex

\def\headdate{14th September 2024}

\section{Steady state viscoelastic fluid flow obtained using Newton-Raphson iteration}
\label{sec:steadyviscoelastic}

In this section some examples will be given that obtain the steady state solution
of the equations for viscoelastic fluid flow using Newton-Raphson iteration. For that, we will make use of the
Jacobian of the nonlinear equations used for Newton-Raphson iteration in implicit time-discretization schemes of
Sec.~\ref{sec:implicit}.

The steady state solution must fulfill the equations Eqs.~\eqref{eq:impve1}--\eqref{eq:impve3} with the time-derivatives $\plderiv{\vek u}t$ and $\plderiv{\ten c}t$ set to zero.
When discretizing the steady state equations in space, we obtain a set of nonlinear equations:
\begin{equation}
  \col F(\col u) = \col 0
\end{equation}
where $\col u$ is the column vector of unknown degrees of freedom. The Jacobian of $\col F(\col u)$ has
already been determined in the time-discretization schemes of
Sec.~\ref{sec:implicit} and can be reused here to perform a Newton-Raphson iteration for finding a steady state solution.
This is indeed possible for the conformation tensor formulation, since $\plderiv{\ten c}t=\col 0$. However for the contravariant deformation or b-tensor formulation (see Sec~\ref{sec:CDTformulation}) we have seen that the $\ten b$ tensor will not be in a steady state even if $\ten c$ is. This means that $\plderiv{\ten b}t\neq\col 0$ and an additional nonlinear term needs to be added (see \cite{JaenssonHulsen2024}) to take that effect into account:
 \begin{equation}
  \col T(\col b) + \col F(\col u) = \col 0
\end{equation}
where $\col b$ are the b-tensor unknowns.
Hence the Jacobian consists of the sum of the Jacobion of $\col T(\col b)$ and the Jacobion  of $\col F(\col u)$.

\subsection{Planar Couette flow for a Weissenberg number (Wi) of 10}
\label{sec:couette_steady_implicit}

A Couette flow example \texttt{couette27.f90} is available that first computes the steady state using Newton-Raphson iteration instead of using time stepping to a steady state (such as example \texttt{couette22a.f90} in Sec.~\ref{sec:couette_implicit}). Then the flow is perturbed and solved using fully-implicit time integration. It uses element level rotation reinitialization like example \texttt{couette22a.f90}.

We refer to Sec.~\ref{sec:couette} for a description of the Couette problem.

The example can be obtained by typing at the
command line:
\medskip
\begin{quote}
   \texttt{getexample couette}
\end{quote}

\subsection{Flow around a cylinder between two plates (fully implicit and steady state)}
\label{sec:confcyl2}
The examples can be obtained by typing at the command line:
\begin{quote}
   \texttt{getexample confined\_cylinder\_b}
\end{quote}

The examples describe the flow
around a fixed cylinder positioned between two walls. 

The example \texttt{cylinder7.f90} uses
fully-implicit time discretization for a few steps and the last step uses Newton-Raphson iteration to
find the steady state solution. The example is similar to the example \texttt{cylinder6.f90} of Sec.~\ref{sec:cylinderimplicit}, but now with a last step for finding the steady state by iteration.
The example \texttt{cylinder8.f90} computes a sequence of steady state solution for increasing flow rate (Weissenberg number), where each solution is the starting solution for the next iteration process. Example \texttt{cylinder10.F90} is similar to \texttt{cylinder8.f90}, but reads from an input file:
\begin{quote}
   \texttt{cylinder10 < input\_cylinder10.txt}
\end{quote}
and there is also the option to restart from a previously saved solution. The example \texttt{cylinder10.F90} is used to create the results in \cite{JaenssonHulsen2024}.

Examples \texttt{cylinder\_c7.f90}, \texttt{cylinder\_c8.f90} and \texttt{cylinder\_c10.F90} are also available
and similar to \texttt{cylinder7.f90}, \texttt{cylinder8.f90} and \texttt{cylinder10.F90}, respectively, however the problems are solved using the conformation tensor approach. 

\subsection{Flow around a sphere at the center of a tube}
\label{sec:settsphere}
The examples can be obtained by typing at the command line:
\begin{quote}
   \texttt{getexample settling\_sphere\_b}
\end{quote}

The examples describe the steady flow due to a sphere moving at constant velocity inside a cylinder. 
The problem is solved in the frame of reference of the moving particle, i.e.~the velocity on the
cylinder wall as well as the left and right boundaries is equal but opposite to the particle 
velocity, whereas the particle itself is stationary.

The example \texttt{settling1.f90} solves this problem on a 2D mesh (assuming axisymmetry) whereas
\texttt{settling2.f90} is a full 3D simulation using an angular cutout of the 3D domain. 
For the latter case, local transformations (see Ch.~\ref{ch:transformations}) are applied such that 
the normal component of the velocity and the tangential components of the traction are zero on 
the sides of the domain, which are created by the cutout. Both \texttt{settling1.f90} and \texttt{settling2.f90}
compute a sequence of steady state solutions for increasing particle velocity 
(Weissenberg number), where each solution is the starting solution for the next iteration process. 

The \texttt{meshes} folder contains a 2D and 3D mesh, as well as code to 
generate new meshes using Gmsh (see Sec.~\ref{sec:readgmshmesh}). 
The relevant files are \texttt{generate\_meshp.f90} and \newline
\texttt{generate\_meshp\_3d.f90} for 2D and 2D meshes, respectively. Mesh characteristics
can be changed in these files.

Examples \texttt{settling3.F90} and \texttt{settling4.F90} are similar to \texttt{settling2.f90}, 
but use Metis (see Sec.~\ref{sec:metis5}) for renumbering and Pardiso (see Sec.~\ref{sec:pardiso}) as solver, respectively.
The example \texttt{settling4.F90} is used to create the results in \cite{JaenssonHulsen2024}.
